%=== CHAPTER TWO ===
%% Rewritten 2026-09-18: adds inertial-parameter identification, aerial
%% manipulation and change detection, and restates the gap in eventless terms.

\chapter{Background and Related Work}
\label{ch:literature}
\begin{spacing}{1.5}
\setlength{\parskip}{0.22in}

\section{Model Predictive Control}

MPC repeatedly solves a finite-horizon optimal control problem, applies the first
input and shifts the horizon forward. Its practical value comes from combining
prediction, multivariable coordination and explicit constraints within one
optimisation \cite{rawlings2017mpc,mayne2014mpc}. Linear MPC is mature and
computationally efficient, but a single linear model is poorly suited to the
large attitude changes and thrust-direction coupling of a quadrotor. NMPC retains
the nonlinear dynamics and nonlinear tracking errors directly
\cite{grune2017nmpc}.

The price is computation. Multiple shooting, sequential quadratic programming,
Gauss--Newton Hessian approximations and real-time iteration schemes are widely
used to make NMPC practical \cite{diehl2005rti}. CasADi provides symbolic
differentiation and a high-level optimisation interface
\cite{andersson2019casadi}; acados generates efficient optimal-control solvers
around tailored SQP and QP components \cite{acados,verschueren2022acados}, with
HPIPM \cite{hpipm} as the quadratic subproblem solver. The present work uses full
SQP with merit-function backtracking, because the real-time iteration variant was
observed to diverge silently on this problem.

\section{Quadrotor NMPC}

A quadrotor is differentially flat under common assumptions \cite{mellinger2011},
and many successful controllers exploit position and yaw references with a
cascaded inner attitude loop \cite{lee2010}. NMPC can instead optimise collective
thrust and body torque against the complete nonlinear rigid-body model, and this
formulation becomes advantageous as trajectory curvature and operating envelope
increase \cite{kamel2017linear,foehn2021}.

The controller in this dissertation predicts position, velocity, unit quaternion
and body rate. The attitude residual is the vector part of the relative
quaternion, which avoids treating the four quaternion components as independent
Euclidean coordinates. Input limits represent collective thrust and body-torque
authority. This formulation is nominally standard; the distinguishing issue is
that the equilibrium input and the plant parameters change at payload attachment
and again at release.

\section{Uncertainty and Adaptation in MPC}

Robust MPC protects against bounded uncertainty, while stochastic MPC describes
uncertainty probabilistically. Both can preserve margins without identifying the
current plant, but conservatism grows when the uncertainty set must cover both an
empty and a loaded vehicle. Adaptive MPC instead updates a parameterised model
online: certainty-equivalent schemes use the current estimate directly,
set-membership schemes maintain a feasible parameter set, and dual-control
schemes excite the plant to learn \cite{aswani2013learning,hewing2020learning}.

Learning-based MPC extends this using Gaussian processes, neural networks or
learned residual models. Data-driven MPC for quadrotors can compensate complex
aerodynamic effects \cite{torrente2021datadriven,torrente2021}, and neural
approximations can reduce online computation \cite{salzmann2023neural}; broader
surveys of learning for control under safety constraints are available
\cite{brunke}. Those methods are valuable when the uncertainty is poorly
structured. Payload pickup has stronger prior structure: rigid-body mechanics
already specifies how mass and attachment geometry enter the equations. The
present work therefore estimates a small physical parameter set rather than
learning a general residual map.

A distinct and influential alternative is to decline to identify the payload at
all, and to absorb its effect as a lumped matched disturbance. The
$\mathcal{L}_1$-adaptive augmentation of a quadrotor NMPC \cite{hanover2021} is
the representative of this family, and incremental nonlinear dynamic inversion
\cite{smeur} is its inner-loop counterpart. Such schemes are robust and require
no payload model, but their internal quantity is a force or an acceleration, not
a mass and a moment. They therefore cannot answer the question this dissertation
poses --- whether a payload is still attached --- because that question is not
represented in their state. The comparison is implemented here (an additive
lumped translational and rotational disturbance in the prediction model) but is
not among the frozen results.

\section{State and Parameter Estimation}

Recursive filters such as the extended Kalman filter propagate a local Gaussian
approximation. They are efficient and central to flight-state estimation, but
inequality constraints are not naturally hard constraints and abrupt parameter
changes violate the local linearisation assumption. MHE estimates a window of
past states and disturbances by solving a constrained optimisation; it can impose
bounds such as a physically admissible mass interval and can attribute residuals
using the same nonlinear model as the controller
\cite{rawlings2017mpc,rao2003mhe,rao2003}. Recent work learns the MHE weights
themselves \cite{neuromhe}, which is the principled version of the stage-weight
question that Chapter~\ref{ch:framework} settles negatively for this problem.

A moving window introduces a trade-off that this work encounters directly. A
longer horizon improves noise rejection, but the constraint that makes an
augmented parameter identifiable --- that it is constant across the window ---
becomes less defensible as the window spans more of the trajectory. The
implemented estimator uses 20 stages at \SI{0.1}{\second}.

\section{Inertial-Parameter Identification}

That the rigid-body equations are \emph{linear} in a particular set of inertial
parameters --- mass, first mass moment and the six inertia elements --- is a
classical result in manipulator identification \cite{atkeson1986}. The first mass
moment $\svec = \mP\bm r_{xy}$ used throughout this dissertation is a member of
that set, and the argument of Chapter~\ref{ch:model} is essentially an
application of the classical observation to a case where it has not usually been
applied: an aerial vehicle whose payload can \emph{leave}, so that the factored
parameterisation becomes singular while the product does not.

Payload-aware estimation for multirotors has previously used force--torque
sensing, external observers or explicit mass estimators \cite{tomic2017,
svacha2020, wuest2019, bohm2021}, and slung-load and aerial-manipulation work
treats the payload as a dynamic subsystem of its own \cite{sreenath2013,
thomas2014, ollero2022}. Both families assume that whether a payload is present
is known. That assumption is what this dissertation removes.

\section{Change Detection and Event-Triggered Estimation}

Deciding that a plant has changed, from a residual, is the subject of classical
change-detection theory \cite{basseville1993}, and event-triggered estimation is
an established response to it \cite{heemels2012}. The predecessor of the present
system used exactly this structure: a persistent step in the reconstructed
collective thrust triggered a de-weighting of the pre-event stages of the window.

Chapter~\ref{ch:framework} reports why that structure was retired for this
problem. A collective-thrust residual cannot distinguish a payload change from a
trajectory-mode change, a sustained vertical gust, or an aggressive manoeuvre,
because all of them move the same scalar. Any decision resting on that signal
alone inherits the ambiguity. The deployed design instead requires evidence from
the \emph{rotational} channel --- a collapse of the first moment relative to a
self-formed reference --- which is the one channel that measured cleanly in the
development data. A formal comparison against a CUSUM detector on the same
signals, under a matched false-alarm constraint, is not performed and is listed
among the gaps.

\section{Research Gap}

The literature contains strong results on quadrotor NMPC, on learned residual
dynamics, on inertial-parameter identification and on payload-aware aerial
control. Three things are missing at their intersection, and this dissertation
addresses them.

First, \textbf{a payload parameterisation that survives release.} Existing
payload-adaptive formulations assume attachment and estimate parameters
conditioned on it. None of them is well posed in the limit $\mP\to0$ with a
non-zero assumed lever arm, which is the regime in which every release occurs.

Second, \textbf{a model-update rule that does not consult the event.} Systems
that grasp and release overwhelmingly treat the gripper command, or a contact
sensor, as ground truth for the model. The mismatch between command and physical
detachment is not usually measured, and the failure it produces --- a loaded
vehicle flown with an empty model --- is silent.

Third, \textbf{an honest account of what such a design costs.} Refusing the
command delays confirmation, and the delay is bimodal rather than merely larger.
Estimating an absolute mass from an absolute thrust makes the whole interface
sensitive to a calibration constant. Both are quantified in
Chapter~\ref{ch:evaluation} rather than asserted to be small.

The dissertation consequently treats estimation and control as one architecture
rather than two algorithms. Estimation is judged by whether its outputs are
identifiable and useful for control; control is judged by whether it remains
consistent with the quantities being estimated, and by whether it knows when it
does not know.

\end{spacing}
\newpage
